\section{Related Work}

\paragraph{\textbf{Heuristic heap sizing.}}
Brecht et al.~\cite{brecht:nogc} propose a heuristic based resizing mechanism for expanding the heap in Boehm-Demers-Weiser GC~\cite{boehm:gc}.
They aim to reduce \flexheapGC{} cost while avoiding paging.
The heap grows by different amounts, depending on its current size in relations to a set of predefined
hand-tuned threshold values. %However, the heap size is never reduced.
%
PAMM~\cite{zhang:pamm} exploits repetition in program phase behavior to manage memory adaptively.
However, their technique does not apply to complex real-life big data frameworks.
%
IslaVista~\cite{grz:isla_vista} proposes a heap resizing policy to avoid \flexheapGC{}-induced paging using OS information.
Heap grows linearly as there are not allocation stalls and shrinks aggressively when allocation stalls are detected.
%
Bruno et al.~\cite{bruno:resize} propose a sizing policy enabling JVM to dynamically grow and shrink the heap according to the cloud application requirements.
%
Cook et al.~\cite{cook:auto} and Sim{\~a}o et al.~\cite{simao:adaptability} study how the \flexheapGC{} heap resizing policy affects the performance of the application.
%
Tavakolisomeh et al.~\cite{zgc:heap_resize} propose the heap size adjustment based on the CPU time spent on \flexheapGC{}.
This approach focuses on minimizing CPU overhead for concurrent collectors, ensuring efficient memory use without stalling application performance.
%
In contrast, \flexheapName{} tackles a broader challenge by dynamically partitioning DRAM between the heap and I/O cache, reducing the combined overhead of \flexheapGC{} time and I/O wait time.

\paragraph{\textbf{Explicit model-based heap sizing.}}
Yang et al.~\cite{yang:real_mem} proposed an analytical model to adjust heap size in multi-tenant environments.
The model monitors application memory allocation and footprint, periodically changing heap size to minimize \flexheapGC{} cost and paging.
However, their approach needs OS modifications.
%
Sun et al.~\cite{sun:multi_space_heap} tackle the issue of memory management for many applications operating under a single JVM by proposing segregated, per-application heaplets.
They present an analytical model that modifies heaplet sizes to equilibrate \flexheapGC{} frequency among all applications, so preventing any single application from unduly impacting the system's overall performance.
%
Tay et al.~\cite{tay:page_fault} formulate a page fault equation that correlates the frequency of page faults with heap size and resident heap size, predicated on established \flexheapGC{} and OS rules for a certain program input.
They propose a heap sizing rule that reduces page faults.
However, alterations in input or system load necessitate parameter recalibration, hence constraining its adaptability.
%
Vengerov~\cite{vengerov:gc_analysis} creates a mathematical model to optimize throughput in generational heaps by tuning \flexheapGC{} parameters.
%
Singer et al.~\cite{singer:economics} apply microeconomic supply and demand theory to model \flexheapGC{} behavior, using elasticity to guide heap size expansion.
However, they require from users to set a target elasticity value which is not intuitive.
%
MemBalancer~\cite{kirisame:membalancer} proposes a model-based heap resizing policy for V8 that derives a compositional square-root heap-limit rule.
Its key idea is to coordinate heap allocations across multiple concurrent runtimes under a shared memory budget so as to optimize the trade-off between total GC time and memory usage.
%
Unlike these approaches, \flexheapName{} does not rely on an explicit analytical or predictive model of heap behavior.
Instead, it uses a model-free, feedback-driven controller based on observed changes in GC and I/O cost, which avoids assuming a fixed heap--I/O relationship across workloads and storage environments.

\paragraph{\textbf{Control-theoretic heap sizing.}}
White et al.~\cite{white:pid} suggest a Proportional-Integral-Derivative (PID) controller that monitors the \flexheapGC{} overhead and modifies the heap resize ratio to uphold a user-defined target \flexheapGC{} goal.
%
Storm et al.~\cite{storm:db2} examine autonomic database configuration through control theory to create a self-tuning memory manager that dynamically modifies heap sizes in databases.
Their methodology employs separate heaps for different memory-intensive operations, including SQL cache and buffer pool.
The suggested strategy employs a cost/benefit estimation model to adjust the size of each heap, aiming to equilibrate the cost/benefit metrics among all heaps.
They utilize a multi-input multi-output (MIMO) controller with integral control for tuning, highlighting the advantages of this controller-based approach, which encompasses rapid convergence, swift adaptability, and stable responses to noise.

\paragraph{\textbf{Cross-runtime heap sizing.}}
Prior work~\cite{alonso:advisor,hertz:waste_not,cameron:forseti,wang:elasticMem} has studied dynamic heap sizing in shared environments, adjusting heap sizes across multiple managed processes or VMs to reduce paging or improve overall throughput.
ElasticMem~\cite{wang:elasticMem}, for example, reclaims heap memory from data analytics applications and returns it to the OS for use by other processes.
Our work differs in both scope and abstraction.
Prior systems operate at the JVM abstraction boundary, modeling heap size as a function of throughput without distinguishing the causes of performance degradation.
\flexheapName{}, in contrast, explicitly measures I/O stall time and separates GC displacement from blocked mutator progress.
This cross-layer attribution lies outside the black-box interface assumed by prior multi-runtime allocators, and is important because heap--I/O interactions are mediated by I/O cache dynamics rather than by heap size alone.

\paragraph{\textbf{I/O cache and application coordination.}}
P2Cache~\cite{lee:p2cache} proposes an application-specific I/O cache that enables developers to build a custom I/O cache that aligns with I/O characteristics of their application.
%
Symbiosis~\cite{dai:symbiosis} dynamically configures the size of user-space block caches in persistent key-value stores with the I/O cache to improve their performance, requiring offline profiling.
%
Unlike these systems, \flexheapName{} focuses on the dynamic DRAM division between the JVM heap and I/O cache without application knowledge or offline profiling.
%
Also, \flexheapName{} could be applicable for systems~\cite{khanh:polar,teraheap,kolokasis:teracache,kolokasis:teraheap_toplas,kolokasis:smartsweep} that use I/O cache to extend the JVM heap over slower memory tiers, such as fast storage devices or remote memory.
